% Preamble requirements:
%   \usepackage{amsmath}
%   \usepackage{booktabs}
%   \DeclareMathOperator*{\argmax}{arg\,max}
%   \DeclareMathOperator{\IoU}{IoU}
%
% GENERATED by scripts/ablation_study/write_section.py -- every number here, in the
% prose as well as the tables, is computed from the study's own artefacts.
% Edit the script, not this file.
%
% Verified: compiles in a twocolumn article, no errors, no overfull boxes.

\subsection{Evaluating graph creation}

We ask which design choices for representation and graph creation are best (Q1), and how sensitive the result is to them (Q2), scoring both against ground-truth subtask annotations with a segmentation F1.

For episode $e \in \mathcal{E}$, the segmenter's boundaries partition it into predicted segments $\mathcal{P}_e$ and the annotation gives ground-truth segments $\mathcal{G}_e$, all half-open intervals in seconds. Let $\lvert \cdot \rvert$ be duration, $\#$ cardinality, $M = \sum_e \# \mathcal{P}_e$, $N = \sum_e \# \mathcal{G}_e$, and $t \in (0, 1]$ the IoU threshold.

\begin{align}
  \IoU(p, g) &= \frac{\lvert p \cap g \rvert}{\lvert p \cup g \rvert}
    \label{eq:iou} \\
  a(p) &= \argmax_{g \,\in\, \mathcal{G}_e} \, \lvert p \cap g \rvert
    \label{eq:assign} \\
  r(g) &= \argmax_{p \,\in\, a^{-1}(g)} \, \lvert p \cap g \rvert
    \label{eq:rep} \\
  \mathrm{TP}(t) &= \sum_{e \in \mathcal{E}}
    \# \bigl\{\, g \in \mathcal{G}_e : \IoU\bigl(r(g), g\bigr) \ge t \,\bigr\}
    \label{eq:tp} \\
  F_1(t) &= \frac{2\,\mathrm{TP}(t)}{M + N}
    \label{eq:segf1}
\end{align}

Each prediction is assigned to the ground-truth segment it overlaps most \eqref{eq:assign}, and each $g$ keeps its largest-intersection assignee \eqref{eq:rep}; the rest are duplicate false positives, so $\mathrm{FP} = M - \mathrm{TP}$ and $\mathrm{FN} = N - \mathrm{TP}$. Duplicate suppression caps $\mathrm{TP}$ at $N$ while every surplus prediction enters $M$, so over-segmenting cannot raise the score, and a long prediction has low IoU with each segment it spans, so under-segmenting cannot either. $F_1$ is therefore safe to maximise directly, and is single-peaked in segment density over the grid. Unless stated otherwise every number we report is $F_1(0.5)$: a prediction counts as a hit when it overlaps its ground-truth segment by at least half their union.

% Generated over 8 of 8 tasks: bottles, chess, earbuds, key, legos, surgical, towels, whiteboard.
% Design: 72 configurations per representation per task.

Each task is split into fit, selection and reporting episodes; configurations are chosen on the selection split and every number reported comes from episodes no selection decision touched. We sweep a design that crosses representation, reduction, clusterer, $K$ and smoothing completely, at 72 configurations per representation per task over 8 tasks, and report the mean and standard deviation over tasks. Because equally spaced boundaries are strong on these tasks, we also report a uniform segmenter emitting each configuration's own segment count.

\subsubsection{Q1: which design choices win}

Averaged over tasks, the winning configuration is representation Proprioception, reduction PCA, clusterer $k$-means, $K$ 10. The resulting configurations beat a count-matched uniform segmenter by +0.175 on average, and on every task (min +0.074, max +0.478).

Averages hide how stable an option is, which is what the std column is for. Proprioception has the highest mean at 0.538 but varies by 0.084 across tasks, while Policy conditioning averages 0.502, within 0.036 of it, and varies by only 0.047. Peak accuracy and portability therefore point at different representations, and a study on a single task would have picked whichever suited that task. Policy encoder is last, averaging 0.494: its features are temporally smooth, so at these cluster counts it emits almost no boundaries.

\subsubsection{Q2: how sensitive the score is}

The reduction is the only choice that consistently matters. It explains a mean $\eta^2$ of 0.348 across tasks and never less than 0.034, while no other stage reaches 0.28 on average; representation is the closest at 0.231.

The practical reading is that the method is robust to everything except the reduction, which is the one stage worth spending effort on. Two caveats bound this. An additive model over all 4 stages explains 0.637 of the variance on average (min 0.507, max 0.797), so the remainder lies in combinations of options rather than in any option by itself. And the 8 tasks are one robot and one annotation protocol, so a standard deviation over them describes this benchmark rather than estimating a population.

\begin{table*}[t]
\centering\small
\begin{tabular}{@{}l rr@{}}
\toprule
option & $F_1$ mean & $F_1$ std \\
\midrule
\multicolumn{3}{@{}l@{}}{\textit{representation} \quad $\eta^2 = 0.231$} \\
\quad DINOv3 (frozen) & 0.538 & 0.130 \\
\quad Policy conditioning & 0.502 & 0.047 \\
\quad Policy DiT block & 0.505 & 0.052 \\
\quad Policy encoder & 0.494 & 0.067 \\
\quad Proprioception & 0.538 & 0.084 \\
\quad state_dinov3 & 0.533 & 0.120 \\
\addlinespace
\multicolumn{3}{@{}l@{}}{\textit{reduction} \quad $\eta^2 = 0.348$} \\
\quad UMAP & 0.474 & 0.065 \\
\quad PCA & 0.576 & 0.114 \\
\addlinespace
\multicolumn{3}{@{}l@{}}{\textit{clusterer} \quad $\eta^2 = 0.006$} \\
\quad GMM & 0.545 & 0.073 \\
\quad HSMM & 0.563 & 0.088 \\
\quad $k$-means & 0.569 & 0.101 \\
\addlinespace
\multicolumn{3}{@{}l@{}}{\textit{$K$} \quad $\eta^2 = 0.052$} \\
\quad 4 & 0.477 & 0.087 \\
\quad 6 & 0.525 & 0.088 \\
\quad 8 & 0.544 & 0.095 \\
\quad 10 & 0.563 & 0.111 \\
\quad 12 & 0.562 & 0.109 \\
\quad 16 & 0.560 & 0.098 \\
\bottomrule
\end{tabular}
\caption{Best segmentation $F_1(0.5)$ per option, over 8 tasks on a balanced design of 72 configurations per representation each. Mean and standard deviation are over tasks, so a large std means the option's value is task-specific and a recipe that fixes it will not transfer. $\eta^2$ is the share of grid-wide variance the whole stage explains, averaged over tasks; it is a property of the stage rather than of any one option, so it heads each block, and it ranks them reduction 0.348, representation 0.231, $K$ 0.052. Best configurations score 0.577 on average against 0.402, a margin of +0.175 (min +0.074, max +0.478), for a count-matched uniform segmenter emitting each configuration's own segment count.}
\label{tab:ablation}
\end{table*}

\begin{table*}[t]
\centering\small
\begin{tabular}{@{}ll r@{\hskip 1.2em}r r r@{}}
\toprule
stage & option & mean & $F_1(0.25)$ & $F_1(0.5)$ & $F_1(0.75)$ \\
\midrule
representation & DINOv3 (frozen) & 0.501 & 0.688 & 0.538 & 0.278 \\
 & Policy conditioning & 0.465 & 0.692 & 0.502 & 0.201 \\
 & Policy DiT block & 0.462 & 0.693 & 0.505 & 0.189 \\
 & Policy encoder & 0.457 & 0.691 & 0.494 & 0.185 \\
 & Proprioception & 0.493 & 0.690 & 0.538 & 0.252 \\
 & state_dinov3 & 0.488 & 0.681 & 0.533 & 0.249 \\
reduction & UMAP & 0.438 & 0.660 & 0.474 & 0.180 \\
 & PCA & 0.535 & 0.735 & 0.576 & 0.294 \\
clusterer & GMM & 0.509 & 0.701 & 0.545 & 0.280 \\
 & HSMM & 0.516 & 0.712 & 0.563 & 0.274 \\
 & $k$-means & 0.528 & 0.730 & 0.569 & 0.284 \\
$K$ & 4 & 0.458 & 0.658 & 0.477 & 0.238 \\
 & 6 & 0.493 & 0.690 & 0.525 & 0.264 \\
 & 8 & 0.505 & 0.704 & 0.544 & 0.267 \\
 & 10 & 0.523 & 0.723 & 0.563 & 0.284 \\
 & 12 & 0.521 & 0.716 & 0.562 & 0.284 \\
 & 16 & 0.524 & 0.726 & 0.560 & 0.287 \\
\bottomrule
\end{tabular}
\caption{Segmentation $F_1$ at three IoU thresholds, averaged over 8 tasks. Each cell is that option's best achievable $F_1$ at that threshold: the configuration is selected at the same IoU it is reported at, on the selection split, and scored on the reporting split. Cells are therefore not all the same configuration, which is what makes each one a fair statement about its own option. ``mean'' averages the three thresholds, as COCO's mAP does, and inherits from that the fact that each threshold was optimised separately.}
\label{tab:ablation-strictness}
\end{table*}
